% genere par split_sec9.py -- contenu issu de sec9.tex (audite), reorganise pour le rapport final.
\section{Mécanisme du ralentissement}

\subsection{Identification du mécanisme : traînée non linéaire, mesurée}
\label{subsec:mecanisme}

Deux mécanismes peuvent expliquer le ralentissement, et ils ont des signatures
\emph{opposées}. \textbf{(A) Échantillonnage préférentiel} : la bulle passerait plus de
temps dans du fluide qui descend, et serait ainsi \emph{portée} vers le bas ; sa
signature est une vitesse verticale du fluide vu par la bulle nettement \emph{négative},
du même ordre que la réduction $\Delta U$. \textbf{(B) Traînée non linéaire}
\cite{liu2024} : le fluide ambiant a une vitesse moyenne nulle, mais $C_D(Re)$ étant
non linéaire, les fluctuations de la vitesse de glissement augmentent la traînée
\emph{moyenne} ; sa signature est un fluide ambiant $\simeq0$ mais une forte
\emph{variance} du glissement.

\paragraph{Mesure.} On reprend la démarche de Liu \emph{et al.}~\cite{liu2024}, qui
définissent la vitesse du fluide \og{}vue\fg{} par la bulle comme une moyenne sur une
coquille de liquide entourant l'interface (ils prennent $1{,}5\,d < r < 2\,d$, avec un
contrôle jusqu'à $3\,d$). On échantillonne donc la vitesse verticale du fluide dans un
cône \emph{supérieur} autour de la bulle (\texttt{src/mechanism.c}, snapshots des runs
$\beta=0{,}15$, $0{,}31$, $0{,}38$), à plusieurs rayons croissants, en retranchant
l'écoulement que la bulle \emph{induit} elle-même, calibré sur le run laminaire, où il n'existe aucune
turbulence à échantillonner (il y vaut $+0{,}24$ à $r=3D$, décroissant, et positif comme
attendu en amont d'une bulle montante). La différence proche$-$lointain est indépendante
du repère mobile (le terme $\mathrm{frame\_uz}$ s'y annule). Le tableau ci-dessous oppose
la réduction totale $\Delta U = u_\infty - \overline{W}$ à la part qu'expliquerait
l'échantillonnage préférentiel.

\begin{center}
\renewcommand{\arraystretch}{1.25}
\begin{tabular}{lcccc}
    \toprule
    Run & $\beta$ & $\Delta U$ & signal préférentiel & part de $\Delta U$ \\
    \midrule
    wt05 & $0{,}15$ & $0{,}66$ & $-0{,}01$ & $1\%$ \\
    wt21 & $0{,}31$ & $4{,}71$ & $-0{,}57$ & $12\%$ \\
    wt32 & $0{,}38$ & $3{,}96$ & $-0{,}09$ & $2\%$ \\
    \bottomrule
\end{tabular}
\end{center}

\paragraph{Verdict.} Une fois l'écoulement induit retranché, le fluide vu par la bulle
est \textbf{quasi nul} : l'échantillonnage préférentiel explique au plus $12\%$ de la
réduction, et ce point isolé est vraisemblablement du bruit (la tendance n'est pas
monotone, $\beta=0{,}38$, pourtant plus intense, ne donne que $2\%$). En regard, les
\textbf{fluctuations de la vitesse de glissement} croissent de $8\%$ à $50\%$ de
$\beta=0{,}15$ à $0{,}38$ (fig.~\ref{fig:mecanisme}) : c'est exactement l'amplitude que
requiert la traînée non linéaire. \textbf{Le ralentissement est donc dominé par la
traînée non linéaire}, en accord avec \cite{liu2024} pour une bulle super-Kolmogorov
($St\gg1$), l'échantillonnage préférentiel, réservé aux bulles sous-Kolmogorov
($St\sim1$), restant marginal. On assume les limites de cette mesure : réalisations
uniques, coquille située à l'intérieur de l'échelle intégrale (approximation), trois
points ; mais le signe et l'ordre de grandeur sont robustes.

\begin{figure}[htbp]
    \centering
    \includegraphics[width=\textwidth]{mechanism.png}
    \caption{\textbf{Mécanisme du ralentissement.} À gauche : vitesse du fluide vue par
    la bulle selon le rayon ; le laminaire (gris) donne l'écoulement induit à
    retrancher. Au centre : la réduction totale (rouge) contre la part expliquée par
    l'échantillonnage préférentiel (violet), l'écart étant la traînée non linéaire. À
    droite : les fluctuations de glissement croissent avec $\beta$, condition de la
    traînée non linéaire.}
    \label{fig:mecanisme}
\end{figure}

\paragraph{Cohérence avec la trajectoire.} Ce diagnostic est cohérent avec ce que l'on
observe :
la bulle suit une trajectoire \textbf{oblique quasi-rectiligne} (déviation ${\sim}5$--$7^\circ$
de la verticale, $Ga=70$ étant juste au-dessus du seuil $Ga_c\simeq44$), le long de
laquelle la vitesse d'ascension présente une \textbf{modulation périodique faible} :
sur le plateau, $w$ fluctue avec un écart-type de $0{,}21$, soit $1{,}7\%$ de
$u_\infty$, à la période du lâcher tourbillonnaire ($T\simeq16{,}2$~u.t.,
$St=0{,}080$). C'est une modulation du glissement bien moins spectaculaire que la
spirale du cas $Ga=237$ de Liu et al., mais relevant du même couplage
bulle--sillage.\footnote{Une version antérieure de ce document décrivait ici une
bulle qui \og{}rampe\fg{}. Cette observation est retirée : elle provenait
d'une erreur de rendu. Les films étaient produits par une coupe \emph{horizontale}
($n=\{0,0,1\}$, caméra \texttt{front}), qui place l'axe d'ascension perpendiculaire à
l'écran et rend la montée géométriquement invisible ; seule subsistait à l'image
l'oscillation latérale. Corrigé dans \texttt{src/render\_lab.c} (coupe verticale
$n=\{0,1,0\}$, caméra \texttt{bottom}). La bulle monte en réalité de façon
régulière, à $97{,}6\%$ de $u_\infty$ au minimum.}
Le seuil de non-fragmentation est lui aussi cohérent : nos $We_t\le0{,}34$ sont bien
sous le $We_c=3$ mesuré par Rivière et al. (2021), d'où l'absence de rupture.

\paragraph{Réserve de résolution : levée.} Nos points couvrent
$Fr'\in[0{,}23;0{,}58]$, comblant le bas de la gamme de Liu et al. Une différence de
maillage doit être signalée : leurs DNS sont à ${\sim}68$ mailles par diamètre,
contre ${\sim}17$ pour notre niveau 7. Ce comptage brut est toutefois trompeur, et
l'étude de convergence du \S\ref{subsec:convergence} montre pourquoi : ce qui pilote
la vitesse d'ascension n'est pas le nombre de mailles sur le diamètre, mais la
résolution de la \textbf{couche limite} $\delta\sim D/\sqrt{Re_b}$. Le saut
lvl6$\to$lvl7 ($+12{,}3\%$) correspond exactement au passage d'une couche limite non
résolue ($\delta/\Delta x=0{,}87$) à une couche limite résolue ($1{,}64$) ; au-delà,
doubler encore la résolution ne déplace $u_\infty$ que de $0{,}32\%$. Le niveau 7 est
donc \textbf{convergé pour la grandeur mesurée}, et l'accord avec Liu et al.\ n'est
pas une coïncidence de résolutions différentes.

\paragraph{Leçon de méthode.} C'est le run \emph{laminaire}, dont la seule
raison d'être était de fournir une référence $u_\infty$, qui a révélé
l'erreur numérique initiale. Sans lui, la campagne aurait conclu que \og{}la turbulence supprime la
remontée de la bulle\fg{}, résultat spectaculaire, cohérent avec la littérature
\cite{legendre2025} et faux. Un cas limite dont la réponse est
connue \emph{a priori} (ici : une bulle doit monter indéfiniment en fluide au
repos) est le meilleur détecteur d'erreur numérique.
